\chapter{Results}
This chapter presents the smart tutor system and student synthesis pipeline. It details the metadata and metrics of the generated datasets, as well as the evaluation results for selected knowledge tracing models. The models are tested on LLM-generated, ASSISTments2009, and real app data. The chapter follows the research questions from the Aims and Objectives section, which are restated below for convenience:

\begin{enumerate}
    \item To what extent can knowledge tracing algorithms be effectively evaluated and compared using synthetic student data generated by large language models?
    \item Which knowledge tracing model demonstrates the strongest predictive performance when both trained and tested on LLM-generated synthetic student data within the smart tutor app?
    \item To what degree do evaluation results on synthetic student data align with results on real educational data, such as ASSISTments and real app users?
\end{enumerate}

\section{Smart Tutor App}
The smart tutor app was successfully implemented. The front end (user interface) is written in Unity, and the backend (server logic) is a REST API written in Python. The project is hosted on a personal Raspberry~Pi~5 (a compact, low-cost computer) server and is publicly accessible at \href{https://dissertation.jakubwj.com}{dissertation.jakubwj.com}. The code for this dissertation is accessible at:

\begin{itemize}
    \item \textbf{Dissertation backend and LaTeX} at \href{https://github.com/jakubwj2/dissertation}{github.com/jakubwj2/dissertation}. This repository contains the server implementation, experiment scripts, the dissertation and proposal \LaTeX{} files, and includes the pyKT fork as a submodule.
    \item \textbf{pyKT fork} at \href{https://github.com/jakubwj2/pykt-toolkit}{github.com/jakubwj2/pykt-toolkit}. The original repository can be found at \href{https://github.com/pykt-team/pykt-toolkit}{github.com/pykt-team/pykt-toolkit}.
    \item \textbf{Unity project} at \href{https://github.com/jakubwj2/SmartTutor}{github.com/jakubwj2/SmartTutor}. This includes all Unity assets. Building and running the client requires a local Unity installation; the dissertation submission includes only a WebGL build.
\end{itemize}

The tutor currently includes four question generators: addition, subtraction, multiplication, and division. In this context, a "question generator" is a software component that systematically creates arithmetic practice questions. Each generator takes a maximum operand argument, which specifies the highest number allowed for the numbers in the questions it produces. For the experiments, the maximum operand was set to 9 for all generators, so each operand then ranged from 0 to 9 (except for the second operand in division, which ranged from 1 to 9 to avoid division by zero). With these settings, the system can generate 390 unique questions across the four arithmetic topics. This total could be increased by raising the operand bounds or creating more generator classes.

\section{Student Data Synthesis}
The student generation pipeline produced 500 synthetic users, totalling 37{,}481 logged question responses. This is approximately 75 per student. The exact user and response counts for each LLM-generated dataset, as well as the real and ASSISTments2009 datasets, are shown in Table~\ref{tab:dataset_metadata}. All synthetic users were successfully exported into the pyKT format and used to train the selected knowledge tracing models.

\begin{table}[H]
    \centering
    \caption{Number of problems generated per LLM-Synthesizers.\label{tab:dataset_metadata}}
    \begin{tabular}{+l^c^c^c^c^c}
        \toprule
        \\
        \rowstyle{\bfseries}
        LLM Synthesizer        & DeepSeek-R1 & Llama3.2 & Mistral & Qwen2.5 & Qwen3.5 \\
        \midrule
        Number of Problem Logs & 7735        & 5383     & 8571    & 8457    & 7335    \\
        \bottomrule
    \end{tabular}
\end{table}

The following tables report the area under the ROC curve (ROC-AUC, a measure of a model's ability to distinguish between classes) and accuracy scores for the trained models. ASSISTments2009 is a publicly available dataset of real user data. Although it does not share questions with the synthetic datasets, it serves as a baseline for comparing the relative performance of the knowledge tracing algorithms.

Tables~\ref{tab:trainauc} and~\ref{tab:trainacc} present AUC (Area Under the ROC Curve) and accuracy scores for the training split at the last epoch. Tables~\ref{tab:validauc} and~\ref{tab:validacc} show the same metrics for the validation split. These tables allow comparison of model performance on training and validation data, where AUC measures the ability to distinguish between classes, and accuracy is the proportion of correct predictions. They also help identify potential overfitting, which occurs when a model performs well on training data but poorly on validation data. Across datasets, the largest training-validation score gaps appeared on ASSISTments2009. Several models had higher AUC and accuracy on the training set than on the validation set for this dataset.

\begin{table}[H]
\centering
\caption{Last epoch \textbf{Train AUC} across LLM Synthesizers and Knowledge Tracing Models.\label{tab:trainauc}}
\begin{tabular}{+l^c^c^c^c^c^>{\bfseries}^c}
\toprule
&\multicolumn{6}{c}{\textbf{Knowledge Tracing Algorithms}}\\ 
\rowstyle{\bfseries}
Synthesizer & AKT & DKT & DTransformer & SAKT & SimpleKT & Average \\
\midrule
ASSISTments 2009 & 0.8522 & 0.8652 & 0.9324 & 0.8433 & 0.9130 & 0.8812 \\
\midrule
DeepSeek-R1 & 0.7887 & 0.8101 & 0.9715 & 0.8029 & 0.7892 & 0.8325 \\
Gemma3.1b & 0.8236 & 0.8293 & 0.9888 & 0.8635 & 0.9901 & 0.8991 \\
Llama3.2 & 0.8806 & 0.8996 & 0.9976 & 0.8924 & 0.9215 & 0.9183 \\
Mistral & 0.8305 & 0.8524 & 0.9990 & 0.8873 & 0.8428 & 0.8824 \\
Qwen2.5 & 0.6006 & 0.8775 & 0.8688 & 0.9361 & 0.8825 & 0.8331 \\
\midrule
\rowstyle{\bfseries}
Average & 0.7960 & 0.8557 & 0.9597 & 0.8709 & 0.8899 & 0.8744 \\
\bottomrule
\end{tabular}
\end{table}

\begin{table}[H]
\centering
\caption{Last epoch \textbf{Train Accuracy} across LLM Synthesizers and Knowledge Tracing Models.\label{tab:trainacc}}
\begin{tabular}{+l^c^c^c^c^c^>{\bfseries}^c}
\toprule
&\multicolumn{6}{c}{\textbf{Knowledge Tracing Algorithms}}\\ 
\rowstyle{\bfseries}
Synthesizer & AKT & DKT & DTransformer & SAKT & SimpleKT & Average \\
\midrule
ASSISTments 2009 & 0.8139 & 0.8264 & 0.8650 & 0.8054 & 0.8595 & 0.8340 \\
\midrule
DeepSeek-R1 & 0.8637 & 0.8674 & 0.9810 & 0.8684 & 0.8637 & 0.8888 \\
Llama3.2 & 0.8456 & 0.8688 & 0.9839 & 0.8821 & 0.9759 & 0.9112 \\
Mistral & 0.8407 & 0.8605 & 0.9849 & 0.8687 & 0.8607 & 0.8831 \\
Qwen2.5 & 0.8534 & 0.8608 & 0.9810 & 0.8628 & 0.8601 & 0.8836 \\
Qwen3.5 & 0.8467 & 0.8620 & 0.9829 & 0.8680 & 0.8651 & 0.8849 \\
\midrule
\rowstyle{\bfseries}
Average & 0.8440 & 0.8576 & 0.9631 & 0.8592 & 0.8808 & 0.8810 \\
\bottomrule
\end{tabular}
\end{table}


\begin{table}[H]
\caption{Last epoch \textbf{Validation AUC} across LLM Synthesizers and Knowledge Tracing Models.\label{tab:validauc}}
\begin{tabular}{+l^c^c^c^c^c^>{\bfseries}^c}
\toprule
\rowstyle{\bfseries}
Synthesizer & AKT & DKT & DTransformer & SAKT & SimpleKT & Average \\
\midrule
Assist 2009 & 0.9620 & 0.9628 & 0.9623 & 0.9007 & 0.9582 & 0.9492 \\
DeepSeek-R1 & 0.8962 & 0.8979 & 0.9905 & 0.8986 & 0.8901 & 0.9147 \\
Llama3.2 & 0.8824 & 0.8838 & 0.9891 & 0.8890 & 0.8900 & 0.9069 \\
Mistral & 0.8844 & 0.8768 & 0.9962 & 0.8884 & 0.8997 & 0.9091 \\
Qwen2.5 & 0.8624 & 0.8603 & 0.9929 & 0.8708 & 0.8592 & 0.8891 \\
Qwen3.5 & 0.8861 & 0.8956 & 0.9920 & 0.8916 & 0.8832 & 0.9097 \\
\rowstyle{\bfseries}
Average & 0.8956 & 0.8962 & 0.9872 & 0.8898 & 0.8967 & 0.9131 \\
\bottomrule
\end{tabular}
\end{table}

\begin{table}[H]
\caption{Last epoch \textbf{Validation Accuracy} across LLM Synthesizers and Knowledge Tracing Models.\label{tab:validacc}}
\begin{tabular}{+l^c^c^c^c^c^>{\bfseries}^c}
\toprule
\rowstyle{\bfseries}
Synthesizer & AKT & DKT & DTransformer & SAKT & SimpleKT & Average \\
\midrule
Assist 2009 & 0.8864 & 0.8864 & 0.8869 & 0.8189 & 0.8816 & 0.8720 \\
DeepSeek-R1 & 0.8608 & 0.8675 & 0.9831 & 0.8658 & 0.8565 & 0.8868 \\
Llama3.2 & 0.8699 & 0.8603 & 0.9665 & 0.8689 & 0.8679 & 0.8867 \\
Mistral & 0.8586 & 0.8580 & 0.9823 & 0.8509 & 0.8586 & 0.8817 \\
Qwen2.5 & 0.8443 & 0.8708 & 0.9721 & 0.8680 & 0.8708 & 0.8852 \\
Qwen3.5 & 0.8665 & 0.8665 & 0.9735 & 0.8665 & 0.8525 & 0.8851 \\
\rowstyle{\bfseries}
Average & 0.8644 & 0.8683 & 0.9607 & 0.8565 & 0.8646 & 0.8829 \\
\bottomrule
\end{tabular}
\end{table}


Tables~\ref{tab:testauc} and~\ref{tab:testacc} show ROC-AUC and accuracy scores on the held-out test split. On the synthetic datasets, DTransformer achieved the highest test AUC under most conditions and consistently outperformed other models. On the ASSISTments2009 dataset, the five models had similar AUC and accuracy values, and no single model dominated both metrics.

\begin{table}[H]
\centering
\caption{\textbf{Test Accuracy} across LLM Synthesizers and Knowledge Tracing Models.\label{tab:testacc}}
\begin{tabular}{+l^c^c^c^c^c^>{\bfseries}^c}
\toprule
\rowstyle{\bfseries}
Synthesizer & AKT & DKT & DTransformer & SAKT & SimpleKT & Average \\
\midrule
ASSISTments 2009 & 0.7829 & 0.7821 & 0.7728 & 0.7347 & 0.7602 & 0.7666 \\
\midrule
DeepSeek-R1 & 0.8593 & 0.8650 & 0.9794 & 0.8798 & 0.8728 & 0.8913 \\
Llama3.2 & 0.8707 & 0.8683 & 0.9647 & 0.8707 & 0.8664 & 0.8882 \\
Mistral & 0.8651 & 0.8651 & 0.9712 & 0.8437 & 0.8651 & 0.8820 \\
Qwen2.5 & 0.8357 & 0.8578 & 0.9801 & 0.8556 & 0.8578 & 0.8774 \\
Qwen3.5 & 0.8599 & 0.8599 & 0.9807 & 0.8592 & 0.8552 & 0.8830 \\
\midrule
\rowstyle{\bfseries}
Average & 0.8456 & 0.8497 & 0.9415 & 0.8406 & 0.8462 & 0.8647 \\
\bottomrule
\end{tabular}
\end{table}

\begin{table}[H]
    \centering
    \caption{\textbf{Test AUC} across LLM Synthesizers and Knowledge Tracing Models.\label{tab:testauc}}
    \begin{tabular}{+l^c^c^c^c^c^>{\bfseries}^c}
        \toprule
                         & \multicolumn{6}{c}{\textbf{Knowledge Tracing Algorithms}}                                                       \\
        \rowstyle{\bfseries}
        Synthesizer      & AKT                                                       & DKT    & DTransformer & SAKT   & SimpleKT & Average \\
        \midrule
        ASSISTments 2009 & 0.7851                                                    & 0.7949 & 0.7721       & 0.6667 & 0.7467   & 0.7531  \\
        \midrule
        DeepSeek-R1      & 0.7801                                                    & 0.7742 & 0.9317       & 0.7709 & 0.7706   & 0.8055  \\
        Gemma3.1b        & 0.8277                                                    & 0.8247 & 0.9722       & 0.8270 & 0.9413   & 0.8786  \\
        LLama3.2         & 0.8851                                                    & 0.8920 & 0.9907       & 0.8933 & 0.9207   & 0.9163  \\
        Mistral          & 0.8186                                                    & 0.8371 & 0.9630       & 0.8253 & 0.8244   & 0.8537  \\
        Qwen2.5          & 0.8137                                                    & 0.8541 & 0.8630       & 0.8691 & 0.8677   & 0.8535  \\
        \midrule
        \rowstyle{\bfseries}
        Average          & 0.8184                                                    & 0.8295 & 0.9154       & 0.8087 & 0.8452   & 0.8435  \\
        \bottomrule
    \end{tabular}
\end{table}


\subsection{Observed Behaviour}
Analysis of item-level correctness revealed systematic patterns in the LLM-generated responses. The generated addition and multiplication responses were nearly all correct. For subtraction, approximately half of the responses were incorrect. Inspection of the incorrect subtraction responses revealed that they were exclusively cases in which the correct answer was negative. For division, overall correctness was lower than for addition and multiplication, and higher than for subtraction, with most incorrect responses occurring on questions with correct answers that were non-integer rational numbers.

\begin{table}[H]
    \centering
    \caption{Incorrectly answered across LLM-Synthesizers and Skills.\label{tab:incorrect_answers}}
    \begin{tabular}{+l^c^c^c^c^c}
        \toprule
                       & \multicolumn{5}{c}{\textbf{LLM Synthesizer}}                                          \\
        \rowstyle{\bfseries}
        Skill          & DeepSeek-R1                                  & Llama3.2 & Mistral & Qwen2.5 & Qwen3.5 \\
        \midrule
        Addition       & 1                                            & 1        & 1       & 1       & 1       \\
        Division       & 0.897                                        & 0.9039   & 0.8979  & 0.8885  & 0.9004  \\
        Multiplication & 0.9995                                       & 0.9985   & 0.9995  & 1       & 1       \\
        Subtraction    & 0.5648                                       & 0.556    & 0.5539  & 0.5601  & 0.5563  \\
        \bottomrule
    \end{tabular}
\end{table}